\section{Online Backup Selection under Churn}
\label{sec:online}

Every method in \S\ref{sec:methods} solves a \emph{one-shot} problem: given a
snapshot of host state, choose an assignment. That formulation is correct only if
host state is static. It is not. Hosts degrade and recover, hard failures occur,
and primaries are scaled in and out. The assignment that is best at $t=0$ is not
the assignment that is best at $t=100$.

The obvious repair---re-solve the snapshot problem whenever state changes---is
wrong, and wrong in a way that is worth stating precisely, because it is what the
Hungarian \textsc{optimal} baseline of \S\ref{sec:optimal} would do if it
were deployed rather than evaluated offline. Moving a primary's backup is not
free: the new backup holds no replica of the primary's state, so a migration costs
a full state transfer. A policy that tracks the instantaneous optimum pays that
cost on every fluctuation, including fluctuations that reverse a round later.

This section formulates the online problem, proves that snapshot-chasing has
unbounded competitive ratio, gives an algorithm with a certified bound on what it
can spend on migration, and locates the problem in a complexity class with known
lower bounds.

\subsection{Problem Formulation}
\label{sec:online:model}

Let $V$ be the hosts, $P \subseteq V$ the primaries, and $k = |V| - 1$ the number
of feasible backups for a given primary. Time proceeds in rounds
$t = 1, \dots, T$. At round $t$, assigning primary $p$ to backup $b$ incurs a
\emph{service cost} $c_t(p,b) \ge 0$, which we instantiate with exactly the
weighted feature score of \S\ref{sec:heuristic}, negated and shifted to be
non-negative, so that the online objective is the online version of the paper's
objective rather than a new one.

A policy maintains an assignment $\sigma_t : P \to V$. Writing
$\ell_t(b) = |\{p : \sigma_t(p) = b\}|$ for the load on $b$, the cost of round $t$ is
\begin{equation}
\begin{split}
\mathrm{cost}_t(\sigma_t) = \sum_{p \in P} \Big[\;
  & c_t(p, \sigma_t(p))
  + \lambda\big(\ell_t(\sigma_t(p)) - 1\big) \\
  & + M \cdot \mathbb{1}[\sigma_t(p) \neq \sigma_{t-1}(p)] \Big],
\end{split}
\label{eq:roundcost}
\end{equation}
where $\lambda$ charges congestion when primaries share a backup and $M$ is the
migration cost, expressed in units of one round of service cost. A migration
decided during round $t$ takes effect at round $t+1$: a primary cannot
retroactively have had its state on the new backup. The total cost is
$\sum_{t=1}^{T}\mathrm{cost}_t(\sigma_t)$, and an online policy must choose
$\sigma_t$ knowing $c_s(\cdot,\cdot)$ only for $s \le t$.

An algorithm $A$ is $\rho$-competitive if $\mathrm{cost}(A) \le \rho \cdot
\mathrm{cost}(\mathrm{OPT}) + O(1)$ on every instance, where $\mathrm{OPT}$ is the
minimum-cost schedule chosen with full knowledge of the trace.

Two quantities make the results below readable. The \emph{floor}
$F = \sum_t \sum_p \min_b c_t(p,b)$ is the cost of placing every primary on its
instantaneously cheapest backup with congestion and migration free; no policy,
including $\mathrm{OPT}$, can pay less. The \emph{regret} of a policy that holds
$p$ on $b$ at round $t$ is $c_t(p,b) - \min_{b'} c_t(p,b')$, the excess it is
paying for not being on the best host. Since $F$ is typically the majority of
total cost and no selection policy influences it, we report both the total-cost
ratio and the \emph{excess ratio} $(\mathrm{cost}(A) - F)/(\mathrm{cost}(\mathrm{OPT}) - F)$,
which isolates the part of the bill that selection actually controls.

\subsection{Snapshot-Chasing Is Not Competitive}

\begin{theorem}
\label{thm:snapshot}
Let \textsc{snapshot} be the policy that plays a cost-minimizing assignment for
round $t$ at every $t$. For every $M > 0$ and every $\rho$, there is an instance on
which $\mathrm{cost}(\textsc{snapshot}) > \rho \cdot \mathrm{cost}(\mathrm{OPT})$.
\textsc{snapshot} is therefore not $\rho$-competitive for any $\rho$.
\end{theorem}

\begin{proof}
Take $|P| = 1$, two feasible backups $b_1, b_2$, and $\lambda = 0$. Let
$c_t(p,b_1) = 0$ and $c_t(p,b_2) = \varepsilon$ for odd $t$, and the reverse for
even $t$. The instantaneous minimizer alternates every round, so \textsc{snapshot}
migrates at every round after the first and pays at least $M(T-1)$. The static
schedule that holds $b_1$ throughout pays $\varepsilon \lceil T/2 \rceil$, so
$\mathrm{cost}(\mathrm{OPT}) \le \varepsilon \lceil T/2 \rceil$ and the ratio is at
least $2M(T-1)/(\varepsilon T) \to 2M/\varepsilon$. Choosing
$\varepsilon < 2M/\rho$ gives the claim.
\end{proof}

The construction needs no churn in the usual sense: it is driven entirely by
fluctuations too small to be worth acting on. The failure mode is not that
\textsc{snapshot} chooses badly---it chooses the best host every single round---but
that it has no mechanism to decline a move whose benefit does not repay its cost.
\S\ref{sec:online:results} measures this on the feature-based instances, where
\textsc{snapshot} spends over 80\% of its total cost on migration.

\subsection{A Certified Migration Budget}

\textsc{hysteresis}($\alpha$) accumulates each primary's regret since its last
migration and moves to the current best only once that regret reaches $\alpha M$,
at which point the counter resets. The threshold is expressed in units of $M$, so
the policy needs no tuning when the state size or the link bandwidth changes.

\begin{proposition}
\label{prop:budget}
On every instance, \textsc{hysteresis}($\alpha$) spends at most a $1/\alpha$
fraction of its own service-plus-congestion cost on migration.
\end{proposition}

\begin{proof}
Each migration is preceded by at least $\alpha M$ of accumulated regret, and the
counter is reset on each migration, so regret is never counted twice. After $N$
migrations the policy has therefore accumulated at least $\alpha M N$ regret, giving
$MN \le R/\alpha$ where $R$ is its total regret. The regret charged at round $t$ for
primary $p$ held on $b$ is $c_t(p,b) - \min_{b'}c_t(p,b') \le c_t(p,b)$, because
costs are non-negative, and because migrations take effect only at $t+1$ the policy
actually pays $c_t(p,b)$ in that round. Summing, $R$ is at most the policy's total
service-plus-congestion cost $S$, so the migration spend $MN \le S/\alpha$.
\end{proof}

This is the guarantee Theorem~\ref{thm:snapshot} shows \textsc{snapshot} lacks.
It is a bound on \emph{waste} rather than on competitiveness: it says the policy
cannot be made to thrash, whatever the adversary does, and it holds for every
instance rather than in expectation. The parameter $\alpha$ buys that bound
directly---at $\alpha = 2$ no more than a third of the bill can be migration---and
\S\ref{sec:online:results} reports how much of the budget is used in practice.

\subsection{What Is Not Achievable}

A natural question is whether some cleverer online policy could be
$O(1)$-competitive. It cannot, and the reason is structural rather than specific to
our cost model.

\begin{proposition}
\label{prop:lower}
With $|P| = 1$ and $\lambda = 0$, online backup selection is exactly the metrical
task system (MTS) problem on the uniform metric over $k$ states with switching cost
$M$. Consequently no deterministic online algorithm is better than
$(2k-1)$-competitive, and no randomized algorithm is better than
$\Omega(\log k / \log\log k)$-competitive against an oblivious adversary.
\end{proposition}

\begin{proof}
The states are the $k$ feasible backups; the task vector at round $t$ is
$(c_t(p,b))_{b}$; the distance between any two distinct states is $M$, which is a
uniform metric. The correspondence is cost-preserving in both directions, so the
deterministic lower bound of Borodin, Linial and Saks~\cite{borodin1992mts} and the
randomized lower bound of Karloff et al.~\cite{karloff1994mts} transfer.
\end{proof}

Two consequences matter for practice. First, the competitive ratio necessarily
degrades with the number of candidate backups, so a constant-factor guarantee is
not a reasonable design target at the scales NFV orchestrators operate at; a
budget guarantee of the kind in Proposition~\ref{prop:budget} is. Second, the
reduction means the $(2k-1)$-competitive deterministic work-function algorithm and
the $O(\log^2 k)$ randomized algorithms for uniform MTS are \emph{available} for
this problem. We do not deploy them: both maintain state proportional to the
candidate set and are substantially more complex than a per-primary counter, and
\S\ref{sec:online:results} shows the counter is already close to the offline
optimum on realistic traces. We note the connection because it tells an
implementer where the ceiling is.
